\chapter*{Avant de commencer}
\addcontentsline{toc}{chapter}{Avant de commencer}

\section*{Comment ce logiciel a été fait}
Nous avons développé \soft avec l'aide d'un modèle d'intelligence artificielle, Claude, d'Anthropic, utilisé comme
outil de programmation. Nous préférons le dire clairement, une fois, ici. Le partage des rôles a été simple : nous
avons décidé de ce que devait faire le logiciel, de ce qu'il devait montrer et de la façon de le montrer ; nous
avons testé chaque version sur notre matériel, avec nos propres musiques, et refusé ce qui ne nous convenait pas.
L'outil a écrit le code, les tests et les rendus de contrôle à partir de nos consignes. Aucune version ne vous
parvient sans avoir été compilée, testée automatiquement, puis essayée sur un vrai téléphone.

Le travail décrit dans ce manuel s'étend du 2 octobre au 24 novembre 2026.

\section*{Ce que vous trouverez ici}
Ce manuel vous explique comment utiliser \soft, mais aussi ce qui se passe dans chaque puce sonore et comment le
lire à l'écran. Si vous voulez seulement écouter vos musiques, les deux premières parties vous suffisent. Si vous
composez, si vous étudiez le matériel ou si vous voulez comprendre pourquoi un morceau sonne comme il sonne, la
suite est pour vous.

Quelques conventions : les registres sont écrits en hexadécimal avec le préfixe \$, les notes en notation anglaise
(le la du diapason est A4, à 440~Hz) et les niveaux en décibels par rapport au maximum de la puce. Les captures
sont des rendus de l'application sur un écran de 1080 pixels de large.
